\documentclass[10pt]{article}
\usepackage[margin=.62in]{geometry}
\usepackage[T1]{fontenc}\usepackage{lmodern,microtype}
\usepackage{amsmath,amssymb,booktabs,enumitem}
\setlist{nosep,leftmargin=1.45em}
\pagestyle{plain}
\emergencystretch=2em
\newcommand{\diag}{\operatorname{diag}}
\begin{document}
\begin{center}
{\Large\bf Proof skeleton}\par
{\large Exact Diffusion Design for Maximally Collective Stable Turing Patterns}\\
{\normalsize in Binary-Complex Mass-Action Networks}\\
{\small Topology-Wide All-Spectrum Localization and
Exponent-Optimal Stationary Heterogeneity Trade-Offs}
\end{center}

\section*{1. Reaction algebra and topology-wide localization}
Construct $Y_m$ and $\Gamma_m$ directly from the indexed reactions.
Reaction-by-reaction verification gives $c^T\Gamma_m=0$ for
$c=(0,4,\ldots,4,2,1)^T$.  Solving the balance equations gives exactly
\[
\ker\Gamma_m=\operatorname{span}\{(\mathbf1_m,0,0),(0_m,1,1)\}.
\]
Because $\Gamma_m$ has $m+2$ columns, rank--nullity gives
$\operatorname{rank}\Gamma_m=m$.  Independently, the maximal minor on rows
$X_1,X_3,\ldots,X_m,Z$ and columns
$R_2,\ldots,R_{m-2},R_a,R_b,R_+$ has determinant $4(-1)^m$.
Therefore every positive flux is $(a\mathbf1_m,b,b)$.  Differentiation at the
unit equilibrium gives $A_m(a,b)$; at a general equilibrium
$x^*=H^{-1}\mathbf1$, the Jacobian is $A_m(a,b)H$.  Conversely, the choice
\[
 k_r=\frac{v_r}{(x^*)^{y_r}}
\]
realizes every $a,b,H>0$.

For a nonempty principal set with fewer than $m$ vertices, Frobenius decomposition reduces every nontrivial SCC to: the $X_1,\ldots,X_{m-1}$ long cycle; the $X_2,\ldots,X_m$ long cycle; a principal boundary-triad block; or a negative singleton. At $b=2a$ one directed edge disappears and only refines the decomposition. The two cycle characteristic equations are excluded from $\Re\lambda\ge0$ by strict product inequalities. The full triad has positive cubic coefficients and $c_1c_2-c_3$ is a sum of fourteen positive monomials. Every lower-order principal matrix is thus Hurwitz. Sparse expansion of the $X_1,\ldots,X_m$ block gives the negative signed determinant $-2a^{m-1}b\prod h_i$, forcing a positive real eigenvalue.

For $m=3$, the two long cycles are $\{X_1,X_2\}$ and $\{X_2,X_3\}$,
and every other nontrivial SCC is a principal subgraph of
$\{X_1,X_3,Z\}$.  Thus the chain argument is applied only for $m\ge4$.

\paragraph{Audit object.} Independently generate every induced SCC for $m=3,\ldots,8$ at generic $b$ and at $b=2a$, while the proof uses the all-$m$ graph classification rather than enumeration.

\newpage
\section*{2. Principal-minor diffusion-ray theorem and exact network law}
For $a_I=(-1)^{|I|}\det J_{I,I}$ and $D\succ0$, multilinearity gives
\[
\det(sD-J)=\sum_Ia_I\prod_{j\notin I}sd_j
=s(\beta_1+\beta_2s+\cdots+\beta_ns^{n-1}).
\]
Assume directly that $a_I>0$ for $|I|\le n-2$ and that
$\sum_{|I|=n-1}a_I>0$; a simple zero and otherwise stable spectrum supplies
the latter coefficient condition in the application.  Hence $\beta_k>0$
for $k\ge2$, and the bracket is strictly increasing on $s>0$, proving the
unique simple threshold iff $\beta_1<0$.

For the stronger positive-real band, define
\[
\chi_s(\lambda)=\det(\lambda I+sD-J)
=\sum_Ia_I\prod_{j\notin I}(\lambda+sd_j).
\]
Every derivative contribution from $|I|\le n-2$ is nonnegative. The total
order-$(n-1)$ derivative is $\sum_{|I|=n-1}a_I>0$, the coefficient of
$\lambda$ in $\det(\lambda I-J)$. Hence $\chi_s'(\lambda)>0$ for $\lambda\ge0$. The sign of $\chi_s(0)$ therefore gives a positive real eigenvalue exactly for $0<s<s_*$.
At threshold, $\chi_{s_*}(0)=0$ and $\chi_{s_*}'(0)>0$; this, rather than
scalar simplicity of the generalized-pencil root alone, proves ordinary
algebraic simplicity of zero for $J-s_*D$.

For $A_m(a,b)H$, compute every order-$m$ omission minor. Omitting $Z$ gives weight $-2$; omitting an interior $X_j$ gives $+16$; omitting $X_1$ or $X_m$ gives zero. Thus
\[
\beta_1=2a^{m-1}b\prod_{i=1}^mh_i
\left(8h_Z\sum_{j=2}^{m-1}\frac{d_j}{h_j}-d_Z\right).
\]
This proves the exact criterion. The same sign follows from first-order perturbation of the conservation eigenvalue. For fixed $H$, minimizing $d_{\max}/d_{\min}$ under the strict inequality gives the nonattained infimum $T(H)=8h_Z\sum_{j=2}^{m-1} h_j^{-1}$. At $H=I$, $T=8(m-2)$, and $h_Z/h_j\ge1/\chi_H$ gives $\chi_D\chi_H>8(m-2)$.

\newpage
\section*{3. Improved unit-equilibrium primary bifurcation}
Set $a=b=1$, $H=I$, $K_i=91m-181-i$, and use
\[
d_1=23/63,\quad d_i=K_i^{-1},\quad d_m=1/7,\quad d_Z=16/45.
\]
An explicit right vector satisfies $(A_m-D_m)r=0$ and an explicit left vector satisfies $(A_m-D_m)^T\ell=0$. Exact harmonic-sum formulas prove $\ell^Tr<0$ and $\ell^TD_mr<0$, hence transversality $\eta_m>0$.

Homogeneous stability follows from
\[
\det(\lambda I-A_m)=\lambda q_m(\lambda),\qquad
\lambda q_m=(1+\lambda)^{m-3}P(\lambda)-R(\lambda),
\]
where $P=\lambda^4+12\lambda^3+42\lambda^2+47\lambda+16$ and $R=5\lambda^2+33\lambda+16$. The exact difference $|1+\lambda|^2|P|^2-|R|^2$ expands into 35 nonnegative monomials in $x=\Re\lambda$ and $y^2=(\Im\lambda)^2$, with equality only at the origin.
For $m=3$, $q_3=(\lambda+7)(\lambda^2+5\lambda+2)$ is Hurwitz.  For
$m\ge4$, $|1+\lambda|\ge1$ makes the 35-term strict inequality incompatible
with $(1+\lambda)^{m-3}P=R$ at any nonzero closed-right-half-plane root.
Since $q_m(0)=16m-34>0$, the conservation zero is simple.

For spatial damping $tD_m$, set
$g_1=\lambda+2+(23/63)t$, $g_m=\lambda+5+t/7$,
$g_Z=\lambda+4+(16/45)t$,
$F=g_1g_mg_Z-4g_1-4g_m+g_Z$, and $G=g_Z(4g_1+g_m)-36$.
Chain elimination gives $\det(\lambda I-A_m+tD_m)=Q_mF-G$.
At $(\lambda,t)=(0,1)$, the interior
product telescopes to $91/90$.  With $z=y^2$ and $s=t-1$, the exact source
\[
E_{77}=\left[(91/90)^2+z\right]|F(x+i\sqrt z,1+s)|^2
-|G(x+i\sqrt z,1+s)|^2
\]
has 77 strictly positive coefficients and no constant term.  Thus equality
occurs only at $(x,z,s)=(0,0,0)$. Hence the first mode has one simple zero,
all other first-mode eigenvalues are stable, and all $t>1$ modes are Hurwitz.

The exact unit contrast is
\[
\chi_D=\frac{23(91m-183)}{63},
\]
which improves the superseded high-contrast profile and is within a dimension-independent factor of the sharp linear threshold.

\newpage
\section*{4. Conservation-compatible normal form and stable branch}
The mass-action field is quadratic:
\[
f(\mathbf1+u)=A_mu+\tfrac12B_m(u,u),\qquad c^TB_m(u,v)=0.
\]
On the fixed integrated-mass class, the critical mode is $r\cos\xi$. Since $\cos^2\xi=(1+\cos2\xi)/2$, solve
\[
A_mw_0=-\tfrac14B(r,r),\quad c^Tw_0=0,
\qquad (A_m-4D_m)w_2=-\tfrac14B(r,r).
\]
The first is an augmented nonsingular solve; the second is invertible by mode
isolation.  The interior $w_2$ recurrence telescopes through four consecutive
$K_i$ factors, leaving an explicit $4\times4$ system in the four retained
boundary values.  Its determinant is
$64\mathcal Q_m/[6615(91m-183)(91m-181)(91m-180)]>0$.

The fixed-mass sectorial semiflow has a one-dimensional center manifold.
The spatial reflection $\xi\mapsto\pi-\xi$ commutes with the Neumann
semiflow, preserves the mass class, and sends $r\cos\xi$ to its negative.
An equivariant center manifold therefore has an odd reduced vector field, and
reflection exchanges the two nonzero branches.
Projection of its reduced vector field onto the adjoint critical mode yields
\[
\dot A=\eta_m\mu A+c_mA^3+O(|A|^5+|A|^3|\mu|+|A||\mu|^2),
\]
\[
c_m=\frac{\ell^T\{B(r,w_0)+\tfrac12B(r,w_2)\}}{\ell^Tr}.
\]
Setting this reduced vector field to zero gives the stationary
Lyapunov--Schmidt equation.
Writing $\mathfrak h_m=\sum_{j=1}^{m-2}K_j^{-1}$, the numerator telescopes to
$R_m+C_m\mathfrak h_m$.  The supplement prints $R_m$, $C_m$, their positive
denominators, and the exact $L_m$ clearing identity.  Shifted
coefficientwise-positive polynomials prove the numerator positive, while
$\ell^Tr<0$; therefore $c_m<0$ for every $m\ge3$.
For a direct audit, put
\[
\mathcal Q_m=589180301m^3-3500015940m^2+6930529579m-4574434500,
\]
\begin{align*}
P_R={}&68605040480814208768m^4-550882186169626030957m^3
+1658612632937449670852m^2\\
&-2219226476204103501323m+1113379274975809565700,\\
P_C={}&652054120726848m^4-5151971981328467m^3
+15265080924982572m^2\\
&-20102347725659113m+9927281930180400.
\end{align*}
Then
\[
R_m=\frac{P_R}{286118780220(8m-17)\mathcal Q_m},\qquad
C_m=-\frac{215P_C}{11645046(8m-17)\mathcal Q_m},
\]
and exact clearing gives
\[
R_m+C_m\frac{m-2}{90m-179}
=\frac{L_m}{286118780220(8m-17)(90m-179)\mathcal Q_m}.
\]
The shifted coefficients of $\mathcal Q_m,P_R,P_C,L_m$ are positive, so
$C_m<0$ and
$R_m+C_m\mathfrak h_m\ge R_m+C_m(m-2)/(90m-179)>0$.

Crandall--Rabinowitz gives the stationary local curve; reflection equivariance
and the odd normal form identify its two positive nonconstant halves for
$\mu>0$. Their center eigenvalue is $-2\eta_m\mu+O(\mu^2)<0$. The complementary spectrum is separated at onset and remains separated under analytic perturbation. Positive Neumann diffusion is sectorial on fixed-mass $L^2$, its fractional $H^1$ domain is a Banach algebra in one dimension, and the quadratic Nemytskii map is smooth. The principle of linearized stability gives local exponential convergence in $H^1$ for sufficiently close nonnegative initial data in the same conservation class.

\newpage
\section*{5. Equilibrium-scaled stable trade-off and audit boundaries}
Let $\nu=m-2$ and
\[
 \kappa_\nu=\begin{cases}1/\sqrt3,&\nu=1,\\
 \sqrt5/2,&\nu\ge2,
 \end{cases}
 \qquad L_0=\frac{\kappa_\nu}{\sqrt\nu},
 \qquad L_1=\frac{90\nu}{90\nu+1}.
\]
The interval is nonempty: for $\nu=1$ the squared inequality is
$8281<24300$, while for $\nu\ge2$ it follows from
$(90\nu+1)^2\le(91\nu)^2<6480\nu^3$.  Set $h_1=h_m=h_Z=1$,
$h_i=K_i/(LK_{i-1})$, and choose physical diffusion
$D^{\rm phys}=\mathsf H_m(L)\Delta$. Under
$\widehat x=\mathsf H_m(L)x$,
\[
\partial_t\widehat x=\mathsf H_m(L)
\{f(\widehat x)+(1-\mu)\Delta\partial_{xx}\widehat x\}.
\]
At $\mu=0$, the effective stationary bracket and right vector are unchanged; the dynamic
left vector is $\widetilde\ell_m(L)=\mathsf H_m(L)^{-1}\ell_m$. The physical
fixed-mass functional is $(\mathsf H_m(L)^{-1}c)^T\widehat x$.
Moreover,
\[
\widetilde\ell_m(L)^T\mathsf H_m(L)\Delta r_m
=\ell_m^T\Delta r_m<0,
\qquad \widetilde\ell_m(L)^Tr_m<0,
\]
so the scaled crossing coefficient $\eta_m(L)$ is positive.

Chain elimination now gives interior factors
\[
g_i=\frac{K_{i-1}}{K_i}(1+L\lambda)+\frac{t-1}{K_i}.
\]
For $t\ge1$, Bernoulli's inequality and $\nu L^2\ge1/3$ reduce spatial
half-plane exclusion to the exact source polynomial
\[
E_{84}=(91/90)^2(1+Ax+z/3)|F(x+i\sqrt z,1+s)|^2
-|G(x+i\sqrt z,1+s)|^2,
\quad A=2\nu L, z=y^2, s=t-1.
\]
Its 84 coefficients are coefficientwise nonnegative in $A>0$, and equality
occurs only at the first-mode zero $(x,z,s)=(0,0,0)$.

At $t=0$, the exact homogeneous determinant, after dividing by
$\det\mathsf H_m(L)$, is
\[
 QF_0-R,\qquad
 Q=\prod_{i=2}^{m-1}\left(1+L\frac{K_{i-1}}{K_i}\lambda\right),
\]
\[
 F_0=\lambda^3+11\lambda^2+31\lambda+16,
 \qquad R=5\lambda^2+33\lambda+16.
\]
Every factor of $Q$ has modulus at least $|1+L\lambda|$ for
$\Re\lambda=x\ge0$.  If $\nu\ge2$, then
\[
 |Q|^2\ge1+2\nu Lx+\nu L^2y^2
 \ge1+Ax+\frac54y^2,\qquad A=2\nu L.
\]
On the certified domain $A\ge\sqrt{5\nu}\ge\sqrt{10}>1/4$.
With $z=y^2$ and $U=A-1/4$, the exact source
\[
E_{22}=\left[1+(U+1/4)x+(5/4)z\right]|F_0(x+i\sqrt z)|^2
-|R(x+i\sqrt z)|^2
\]
has 22 nonzero monomials.  All coefficient polynomials are nonnegative and
equality is possible only at $\lambda=0$.  The coefficient $5/4$ is sharp for
this certificate only: replacing it by $B$ makes the $z=y^2$ coefficient at
$x=0$ equal to $256B-320$.  This does not identify an intrinsic dynamical
boundary.  Also
\[
 (QF_0-R)'(0)=16L\sum_{i=2}^{m-1}K_{i-1}/K_i-2>0,
\]
so the conservation zero is simple.  If $\nu=1$, putting
$\gamma=91L/90$ gives
\[
 \frac{QF_0-R}{\lambda}
 =\gamma\lambda^3+(1+11\gamma)\lambda^2
 +(6+31\gamma)\lambda+(16\gamma-2),
\]
whose coefficients and Routh--Hurwitz difference
$6+99\gamma+325\gamma^2$ are positive because $\gamma>1/8$.

The second harmonic is unchanged, while the zero harmonic changes by
$w_0(L)=w_0^{\rm ref}+\tau_m(L)\rho$. The cubic numerator is
$N_m(L)=N_m^{\rm ref}+\tau_m(L)S_m$.  With
$1/91\le\mathfrak h_m<1/90$,
\[
S_m=-\frac{4(1760850\mathfrak h_m-10253)}{462105},
\qquad -\frac1{10}<S_m<0.
\]
Since $C_m<0$, exact clearing gives
\[
N_m^{\rm ref}-\frac1{100}\ge\frac{P_{\rm ref}(\nu)}{D_{\rm ref}(\nu)},
\]
where
\begin{align*}
P_{\rm ref}={}&3790502986637265684840\nu^5
-974216530468600286489\nu^4-53103567440921218871\nu^3\\
&-576386186827093561\nu^2+3649732858601219\nu+55281268032918,\\
D_{\rm ref}={}&715296950550(8\nu-1)(90\nu+1)\\
&\times(589180301\nu^3+35065866\nu^2+629431\nu+3306).
\end{align*}
$D_{\rm ref}>0$, and the shift $\nu=v+1$ gives six positive coefficients for
$P_{\rm ref}$, proving $N_m^{\rm ref}>1/100$.

For the gauge, set
\begin{align*}
A_\tau={}&1494249120\mathfrak h_mL\nu^2-69786990\mathfrak h_mL\nu
+108738630L\nu^2\\
&+1214388L\nu-8521L-125249670\nu^2+1031940\nu,\\
B_\tau={}&32760\mathfrak h_mL\nu+32760L\nu^2+4L-4095\nu.
\end{align*}
Then $\tau_m=-A_\tau/[15876(8\nu-1)B_\tau]$ and
$B_\tau=4095\nu[8L(\nu+\mathfrak h_m)-1]+4L>0$.  The exact derivatives
are negative on $L\ge1/\sqrt{3\nu}$, so the maximum is at
$(\mathfrak h_m,L)=(1/91,1/\sqrt{3\nu})$.  For $\nu\ge2$ its endpoint sign
is bounded by
\[
P_{\rm up}(\nu)=-\frac{189709065}{2}\nu^3-507201030\nu^2-935658\nu+58481,
\]
with
\[
P_{\rm up}'=-\frac{569127195}{2}\nu^2-1014402060\nu-935658<0,
\quad P_{\rm up}(2)=-2789453215<0.
\]
The case $\nu=1$ uses the exact unreduced endpoint expression and
$\sqrt3<7/4$.  Hence $\tau_m(L)<1/20$.  These bounds hold on the broader
cubic range $L\ge1/\sqrt{3\nu}$ and give
$N_m(L)>1/200$, hence $c_m(L)<0$ throughout the certified interval.

The physical contrasts are
\[
\chi_D=\frac{23}{63}91\nu L,\qquad
\chi_H=\frac{91\nu-1}{91\nu L}.
\]
Here $\chi_H$ measures equilibrium concentration-scale separation, not a
universal measure of biological expense.  Moreover
\[
\chi_D(L)\chi_H(L)=\frac{23(91\nu-1)}{63}
=\frac{23(91m-183)}{63}=\chi_D^{\rm unit}(m).
\]
Thus $L$ reallocates a fixed product and reduces
$\max(\chi_D,\chi_H)$, not the product.
At $L_0$,
\[
 \chi_D=\frac{2093}{63}\kappa_\nu\sqrt\nu,
 \qquad
 \chi_H=\frac{\sqrt\nu}{\kappa_\nu}\frac{91\nu-1}{91\nu},
\]
so both are $\Theta(\sqrt m)$. The topology-specific necessary bound for
stationary crossings implies any such realization must have
$\max(\chi_D,\chi_H)>\sqrt{8\nu}$, so exponent $1/2$ is optimal among
stationary-crossing realizations of $\widehat N_m$. Constants remain open.

\paragraph{Failure tests.} Verify the SCC exhaustion at $b=2a$; the
principal-minor derivative argument; the 22/35/77/84-term equality cases and
the direct $\nu=1$ cubic; the transformed physical diffusion; the fixed-mass
gauge; and the comparison $N_m(L)>1/200$. No finite regression substitutes
for these all-$m$ identities.
\end{document}
